%% toc: P2 10 10.3,10.4
\section{Equivalent substitutions and the two congruence properties of equivalence}
\label{sec:p02-w02-p02-004-equivalent-substitutions}

\paragraph{Purpose.}
Section~\ref{sec:p02-propositional-substitution} proves two separate closure
properties of logical equivalence \(\equiv\) on the formulas of
\(\mathsf{Fm}_{\mathsf{CPL}}\).  Theorem~\ref{thm:p02-uniform-substitution}
applies one substitution to two equivalent formulas.
Theorem~\ref{thm:p02-replacement} keeps a context fixed and changes selected
subformula occurrences.  This section combines the two.  Its result answers
the following question: when do two \emph{different} uniform substitutions
give equivalent instances of the same formula, or of two equivalent formulas?
The answer is that pointwise equivalence of the substituted formulas on the
variables that actually occur is sufficient.  Together with the two earlier
theorems it shows that \(\equiv\) is compatible with every primitive
connective and with every uniform substitution, and that the substitution
action depends only on the equivalence classes of the inserted formulas.

\paragraph{Inputs.}
The language, the valuations \(v\) with their recursive extensions
\(\widehat v\), and the relation \(\equiv\) are those of
Sections~\ref{sec:prop-syntax-semantics-01}
and~\ref{sec:prop-syntax-semantics-02}.  The primitive constructors are
\(\bottom,\neg,\land,\lor,\to\); defined connectives are expanded before any
operation below is applied.  Uniform substitutions \(\upsilon\) and their
instances \(\varphi[\upsilon]\) are as in
Definition~\ref{def:p02-uniform-substitution}.  Besides structural
induction, which Theorem~\ref{thm:prop-unique-readability} licenses, the
proofs use only
Theorems~\ref{thm:p02-uniform-substitution} and~\ref{thm:p02-replacement},
both proved in full in the annex.

\begin{hbdefinition}{def:p02-w02-p02-004-occurring-variables}{Occurring variables; equivalent substitutions}
For a formula \(\varphi\), the set \(\operatorname{Var}(\varphi)\) of
propositional variables occurring in \(\varphi\) is defined by structural
recursion:
\[
\begin{aligned}
 \operatorname{Var}(p_n)&=\{p_n\},&
 \operatorname{Var}(\bottom)&=\varnothing,\\
 \operatorname{Var}(\neg\alpha)&=\operatorname{Var}(\alpha),&
 \operatorname{Var}(\alpha\circ\beta)
   &=\operatorname{Var}(\alpha)\cup\operatorname{Var}(\beta),
\end{aligned}
\]
for \(\circ\in\{\land,\lor,\to\}\).  It is a finite set.
Let \(V\subseteq\mathsf{Prop}\).  Two uniform substitutions \(\upsilon\) and
\(\upsilon'\) are \emph{equivalent on \(V\)} if
\(\upsilon(p)\equiv\upsilon'(p)\) for every \(p\in V\).  Nothing is required
of \(\upsilon(p)\) and \(\upsilon'(p)\) for \(p\notin V\).
\index{uniform substitution!equivalent substitutions}
\end{hbdefinition}

\begin{hbtheorem}{thm:p02-w02-p02-004-equivalent-substitutions}{proved=annex:p02-w02-p02-004-eqsub}{Equivalent substitutions, equivalent instances}
Let \(\varphi\) be a formula and let \(\upsilon,\upsilon'\) be uniform
substitutions that are equivalent on \(\operatorname{Var}(\varphi)\).  Then
\[
 \varphi[\upsilon]\equiv\varphi[\upsilon'].
\]
\end{hbtheorem}

\begin{hbcorollary}{thm:p02-w02-p02-004-substitution-invariance}{proved=annex:p02-w02-p02-004-invar}{Joint substitution invariance of equivalence}
Let \(\varphi\equiv\varphi'\), and let \(\upsilon,\upsilon'\) be uniform
substitutions that are equivalent on
\(\operatorname{Var}(\varphi)\cup\operatorname{Var}(\varphi')\).  Then
\[
 \varphi[\upsilon]\equiv\varphi'[\upsilon'].
\]
In particular, if \(\entails\varphi\) and \(\upsilon,\upsilon'\) are
equivalent on \(\operatorname{Var}(\varphi)\), then
\(\entails\varphi[\upsilon]\) and \(\entails\varphi[\upsilon']\).
\end{hbcorollary}

\paragraph{Procedure and side conditions.}
To show \(\varphi[\upsilon]\equiv\varphi'[\upsilon']\) by the corollary,
check three things.
\begin{enumerate}
  \item \(\varphi\equiv\varphi'\), for instance by a truth table over
        \(\operatorname{Var}(\varphi)\cup\operatorname{Var}(\varphi')\).
  \item \(\upsilon(p)\equiv\upsilon'(p)\) for each variable \(p\) of that
        finite set; variables outside it need no check.
  \item Both substitutions are applied uniformly and simultaneously, as in
        Definition~\ref{def:p02-uniform-substitution}: every occurrence of a
        variable receives the same formula, and inserted formulas are not
        substituted again.
\end{enumerate}
The conclusion is an equivalence between the two instances.  It does not say
that \(\upsilon(p)\) and \(\upsilon'(p)\) are equivalent when the instances
are, and it says nothing about operations that treat different occurrences of
one variable differently.

\paragraph{Consequence for the quotient by equivalence.}
Write \([\varphi]\) for the class of \(\varphi\) under \(\equiv\).  By the
congruence clauses of Theorem~\ref{thm:p02-replacement}, the primitive
constructors induce well-defined operations on classes, for example
\(\neg[\alpha]=[\neg\alpha]\) and
\([\alpha]\land[\beta]=[\alpha\land\beta]\).  By
Corollary~\ref{thm:p02-w02-p02-004-substitution-invariance}, for each
substitution \(\upsilon\) the rule \([\varphi]\mapsto[\varphi[\upsilon]]\)
is well defined, and it is unchanged when every \(\upsilon(p)\) is replaced
by an equivalent formula.  So a substitution acts on classes through the
classes \([\upsilon(p)]\) alone.

\begin{hbexample}{def:p02-w02-p02-004-worked-example}{Two different substitutions into two equivalent formulas}
Let \(P,Q,R\) be distinct propositional variables, let
\(\varphi=(P\to Q)\) and \(\varphi'=(\neg P\lor Q)\), so that
\(\varphi\equiv\varphi'\) as in
Example~\ref{def:p02-substitution-replacement-example}.  Put
\[
 \upsilon(P)=(R\land R),\quad \upsilon(Q)=\neg\neg Q,\qquad
 \upsilon'(P)=R,\quad \upsilon'(Q)=Q,
\]
and let both substitutions fix every other variable.  Here
\(\operatorname{Var}(\varphi)\cup\operatorname{Var}(\varphi')=\{P,Q\}\).

\emph{Side condition 2.}  Under any valuation \(v\), the formula
\((R\land R)\) has value \(\mathbf T\) exactly when \(v(R)=\mathbf T\), so
\((R\land R)\equiv R\).  Applying the negation clause twice,
\(\neg\neg Q\) has value \(\mathbf T\) exactly when \(\neg Q\) has value
\(\mathbf F\), that is, exactly when \(v(Q)=\mathbf T\); so
\(\neg\neg Q\equiv Q\).  Hence \(\upsilon\) and \(\upsilon'\) are equivalent
on \(\{P,Q\}\).

\emph{Conclusion.}  The corollary gives
\[
 \varphi[\upsilon]=\bigl((R\land R)\to\neg\neg Q\bigr)
 \;\equiv\;
 (\neg R\lor Q)=\varphi'[\upsilon'].
\]
The annex proof reaches this through the intermediate formula
\(\varphi'[\upsilon]=(\neg(R\land R)\lor\neg\neg Q)\): the first step is
Theorem~\ref{thm:p02-uniform-substitution} applied to
\(\varphi\equiv\varphi'\), and the second is
Theorem~\ref{thm:p02-w02-p02-004-equivalent-substitutions} applied to
\(\varphi'\).  In the second step the leaf \(P\) contributes
\((R\land R)\equiv R\), the leaf \(Q\) contributes
\(\neg\neg Q\equiv Q\), and the congruence clauses for \(\neg\) and \(\lor\)
assemble them.

\emph{Direct check.}  The left formula has value \(\mathbf F\) exactly when
\((R\land R)\) is true and \(\neg\neg Q\) is false, that is, when
\(v(R)=\mathbf T\) and \(v(Q)=\mathbf F\).  The right formula has value
\(\mathbf F\) exactly when \(\neg R\) and \(Q\) are both false, which is the
same condition.
\end{hbexample}

\begin{hbexample}{def:p02-w02-p02-004-limits}{Where the hypotheses cannot be dropped or reversed}
\emph{The hypothesis is needed at occurring variables.}  Take
\(\varphi=P\), \(\upsilon(P)=Q\) and \(\upsilon'(P)=\neg Q\), with \(P\neq Q\).
Then \(\varphi[\upsilon]=Q\) and \(\varphi[\upsilon']=\neg Q\) take different
values under every valuation, so they are not equivalent.  The only variable
of \(\varphi\) is \(P\), and \(Q\not\equiv\neg Q\).

\emph{The hypothesis is sufficient but not necessary.}  Take
\(\varphi=(P\lor\neg P)\) with the same \(\upsilon,\upsilon'\).  The instances
\((Q\lor\neg Q)\) and \((\neg Q\lor\neg\neg Q)\) are both true under every
valuation, hence equivalent, although \(Q\not\equiv\neg Q\).  Equivalence of
instances therefore does not yield equivalence of the substituted formulas.

\emph{Uniformity is essential for validity.}  The formula \((P\to P)\) is
true under every valuation.  Changing only its left occurrence of \(P\) to
\(Q\) gives \((Q\to P)\), which is false when \(v(Q)=\mathbf T\) and
\(v(P)=\mathbf F\).  This is the selected-occurrence replacement of
Definition~\ref{def:p02-selected-replacement} with the non-equivalent pair
\(P,Q\), not a uniform substitution.  Theorem~\ref{thm:p02-replacement}
requires the replaced and replacing formulas to be equivalent, and
Theorem~\ref{thm:p02-uniform-substitution} requires one formula per
variable at all of its occurrences; this operation meets neither condition.
\end{hbexample}

\paragraph{Relationships.}
\begin{itemize}
  \item \emph{Exercise and its feedback.}  The proof of
        Theorem~\ref{thm:p02-w02-p02-004-equivalent-substitutions} uses
        reflexivity of \(\equiv\) and its congruence clauses for every
        primitive connective, and the corollary adds transitivity.  These are
        exactly the facts that Exercise~\ref{ex:prop-metatheory-09} asks for;
        Convention~\ref{def:p02-w02-p02-004-exercise-correspondence} matches
        the exercise and its feedback sketch with the annex proof of
        Theorem~\ref{thm:p02-replacement}.
  \item \emph{The deductive counterparts in Kleene.}  For the propositional
        calculus as a deductive system, Kleene's Theorem~3 states that
        simultaneous substitution of formulas for distinct proposition
        letters preserves deducibility, and in particular provability
        \cite[\S 25, Theorem~3, p.~109]{Kleene1952}.  His replacement
        theorem, Theorem~6, has the form
        \(A\sim B\vdash C_A\sim C_B\), where \(C_B\) comes from \(C_A\) by
        replacing one specified occurrence of \(A\) by \(B\), and \(\sim\) is
        his equivalence connective; it is proved by induction on the depth of
        that occurrence \cite[\S 26, Theorem~6, p.~116]{Kleene1952}.  Its
        second version drops the restriction to proposition letter formulas
        for \(A\) and \(B\), provided \(C_A\) is built from the specified
        occurrence of \(A\) with the propositional operators only
        \cite[\S 26, p.~117]{Kleene1952}.  These are
        statements about derivability; the results of this section and of
        Section~\ref{sec:p02-propositional-substitution} are statements about
        truth values under all valuations, and their proofs use no
        derivations.
  \item \emph{The semantic notions in Kleene.}  Kleene calls two proposition
        letter formulas \emph{identically equal} when their truth tables have
        the same value column, and he observes that adding letters to the
        list does not change identical truth or identical equality
        \cite[\S 28, pp.~127--128]{Kleene1952}.  The relation \(\equiv\)
        used here quantifies over all valuations of \(\mathsf{Prop}\)
        instead of over the rows of one table.  His proof that the axioms are
        identically true checks each axiom schema by assigning truth values
        directly to the schematic parts \(A,B,C\), and then transfers the
        result to every instance \cite[\S 28, Lemma~12a, p.~128]{Kleene1952}.
        That transfer is the validity clause of
        Theorem~\ref{thm:p02-uniform-substitution}.
  \item \emph{Other substitution results in this handbook.}  The first-order
        substitution lemma, Theorem~\ref{thm:fol-substitution}, concerns
        substitution of a term for an individual variable and compares it with
        a modified variable assignment; its operation and its side conditions
        differ from the ones here.  Lemma~\ref{lem:nd-substitution} inserts
        derivations for open assumptions and is a statement about natural
        deduction derivations, not about formulas.  The normal modal
        Lindenbaum algebra of Definition~\ref{def:p10-modal-lindenbaum} is a
        quotient by provable equivalence rather than by \(\equiv\), and its
        well-definedness is proved there for that relation.
\end{itemize}
\index{logical equivalence!substitution invariance}
